\documentclass[12pt]{article}
\usepackage[utf8]{inputenc}
\usepackage[T1]{fontenc}
\usepackage{hyperref}
\title{Tulai Language Resources for K-12 Curriculum Development}
\author{Compiled by Claude Sonnet 5 (claude-sonnet-5) for Andr\'as Kornai}
\date{}
\begin{document}
\maketitle

\section{Introduction}
Tulai (ISO 639-3 tvi; also called Tule) is a West Chadic language of the Afro-Asiatic family, spoken in southern Bauchi State, Nigeria. One catalogue calls it dormant and notes only ``a handful'' of speakers. Search results turned up almost nothing usable for schools, and this report says so plainly.

\section{Primary Educational Resources}
No primers, textbooks or school programmes were found.

\section{Linguistic and Reference Materials}
\begin{itemize}
\item Endangered Languages catalogue entry: \url{https://endangeredlanguages.com/lang/10453}
\item ThoT database entry (tonal status), Glottocode nucl1693: \url{https://thot.huma-num.fr/db/nucl1693/tonal_status}
\end{itemize}
Tulai should not be confused with Tal (Plateau State) or Tula (Gombe State; Waja-Kam, not Afro-Asiatic).

\section{Religious and Community Resources}
\begin{itemize}
\item Joshua Project, Tulai in Nigeria. The profile is generic and gives no linguistic detail: \url{https://joshuaproject.net/people_groups/22370}
\end{itemize}

\section{Digital and Audio Resources}
None were found.

\section{Curriculum Development Context}
With no first-language speakers reported, any curriculum would be a revitalization effort. It would need primary documentation first, namely recordings and wordlists from any remaining semi-speakers, which the search results did not locate.
\end{document}